% !TeX program = XeLaTeX
% !TeX root = ../../shloka.tex
% --meta--
% deity: Shiva
% composer: Markandeya
% type: Ashtakam
% wiki_title: Chandrashekhara Ashtakam
% --end-meta--
% Text: श्रीगङ्गाकल्पोक्तपूजाविधिः (Idkidu, 2015), pp. 54-55.

\sect{चन्द्रशेखराष्टकम्}

\twolineshloka
{चन्द्रशेखर चन्द्रशेखर चन्द्रशेखर पाहि माम्}
{चन्द्रशेखर चन्द्रशेखर चन्द्रशेखर रक्ष माम्} %॥१॥

\fourlineindentedshloka
{रत्नसानुशरासनं रजताद्रिशृङ्गनिकेतनं}
{शिञ्जिनीकृतपन्नगेश्वरमच्युताननसायकम्}
{क्षिप्रदग्धपुरत्रयं त्रिदिवालयैरभिवन्दितं}
{चन्द्रशेखरमाश्रये मम किं करिष्यति वै यमः} %॥२॥

\fourlineindentedshloka
{पञ्चपादपपुष्पगन्धपदाम्बुजद्वयशोभितं}
{भाललोचनजातपावकदग्धमन्मथविग्रहम्}
{भस्मदिग्धकलेवरं भवनाशनं भवमव्ययं}
{चन्द्रशेखरमाश्रये मम किं करिष्यति वै यमः} %॥३॥

\fourlineindentedshloka
{मत्तवारणमुख्यचर्मकृतोत्तरीयमनोहरं}
{पङ्कजासनपद्मलोचनपूजिताङ्घ्रिसरोरुहम्}
{देवसिन्धुतरङ्गसीकरसिक्तशुभ्रजटाधरं}
{चन्द्रशेखरमाश्रये मम किं करिष्यति वै यमः} %॥४॥

\fourlineindentedshloka
{यक्षराजसखं भगाक्षहरं भुजङ्गविभूषणं}
{शैलराजसुतापरिष्कृतचारुवामकलेवरम्}
{क्ष्वेडनीलगलं परश्वधधारिणं मृगधारिणं}
{चन्द्रशेखरमाश्रये मम किं करिष्यति वै यमः} %॥५॥

\fourlineindentedshloka
{कुण्डलीकृतकुण्डलेश्वरकुण्डलं वृषवाहनं}
{नारदादिमुनीश्वरस्तुतवैभवं भुवनेश्वरम्}
{अन्धकान्धकमाश्रितामरपादपं शमनान्तकं}
{चन्द्रशेखरमाश्रये मम किं करिष्यति वै यमः} %॥६॥

\fourlineindentedshloka
{भेषजं भवरोगिणामखिलापदामपहारिणं}
{दक्षयज्ञविनाशनं त्रिगुणात्मकं त्रिविलोचनम्}
{भुक्तिमुक्तिफलप्रदं सकलाघसङ्घनिबर्हणं}
{चन्द्रशेखरमाश्रये मम किं करिष्यति वै यमः} %॥७॥

\fourlineindentedshloka
{भक्तवत्सलमर्चितं निधिमक्षयं हरिदम्बरं}
{सर्वभूतपतिं परात्परमप्रमेयमनुत्तमम्}
{सोमवारिनभूहुताशनसोमपानिलखाकृतिं}
{चन्द्रशेखरमाश्रये मम किं करिष्यति वै यमः} %॥८॥

\fourlineindentedshloka
{विश्वसृष्टिविधायिनं पुनरेव पालनतत्परं}
{संहरन्तमपि प्रपञ्चमशेषलोकनिवासिनम्}
{क्रीडयन्तमहर्निशं गणनाथयूथसमन्वितं}
{चन्द्रशेखरमाश्रये मम किं करिष्यति वै यमः} %॥९॥

\fourlineindentedshloka
{मृत्युभीतमृकण्डुसूनुकृतस्तवं शिवसन्निधौ}
{यत्र कुत्र च यः पठेन्न हि तस्य मृत्युभयं भवेत्}
{पूर्णमायुररोगितामखिलार्थसम्पदमादरं}
{चन्द्रशेखर एव तस्य ददाति मुक्तिमयत्नतः} %॥१०॥

{॥इति श्री-चन्द्रशेखराष्टकं सम्पूर्णम्॥}
